% !TEX root = RELATORIO_NOVO.tex

\documentclass[12pt,a4paper]{article}

\usepackage[T1]{fontenc}
\usepackage[portuguese]{babel}
\usepackage{geometry}
\usepackage{amsmath}
\usepackage{amssymb}
\usepackage{booktabs}
\usepackage{float}

\geometry{margin=2.5cm}

\title{\textbf{Sistema de Recomendação de Perfumes Nacionais} \\ 
       \large Hibridização de TF-IDF e SVD para Resolução do Cold-Start}
\author{Carlos Victor Albuquerque Oliveira - 232009558 \\ Lucas Sena de Almeida - 190112310}
\date{Universidade de Brasília \\ Departamento de Ciência da Computação \\ Turma 01, 2026/1 \\ Prof. Díbio}

\begin{document}

\maketitle

\begin{abstract}
Este projeto implementa um sistema híbrido de recomendação que combina TF-IDF (filtragem por conteúdo) e SVD (filtragem colaborativa) para resolver o problema de cold-start em uma loja virtual de perfumes nacionais. A abordagem híbrida 70\% TF-IDF + 30\% SVD oferece recomendações relevantes e personalizadas mesmo para novos usuários sem histórico de avaliações. O sistema foi validado em uma matriz de 500 usuários × 50 perfumes com 23,6\% de densidade, demonstrando efetividade na combinação de explicabilidade com descoberta colaborativa.
\end{abstract}

\section{Introdução}

Sistemas de recomendação enfrentam o desafio fundamental do \textbf{cold-start}: recomendações para novos usuários sem histórico de interação. As abordagens tradicionais apresentam limitações inerentes:

\begin{itemize}
    \item \textbf{Filtragem por conteúdo pura:} Recomendações genéricas e óbvias
    \item \textbf{Filtragem colaborativa pura:} Impossível sem dados de usuários similares
    \item \textbf{Hibridização:} Combina força de ambas
\end{itemize}

Neste projeto, desenvolvemos um sistema híbrido para uma loja de perfumes nacionais. A motivação é prática: personalizar recomendações mantendo tanto a explicitabilidade (``você selecionou floral'') quanto a descoberta colaborativa (``usuários como você também adoram'').

\section{Definição do Problema e Dados}

\subsection{Catálogo de Produtos}

O projeto trabalha com um catálogo de \textbf{50 perfumes nacionais} de 6 marcas (O Boticário, Natura, Eudora, Avon, Mahogany, Granado, Phebo). Cada perfume possui três características principais:

\begin{enumerate}
    \item \textbf{Família olfativa} (23 tipos): floral, amadeirado, cítrico, oriental, gourmand, aquático, aromático, etc.
    \item \textbf{Ocasião de uso}: categorias diurno/noturno, formal/casual, verão/inverno (ex: ``diurno casual'')
    \item \textbf{Gênero}: masculino, feminino ou unissex
\end{enumerate}

Adicionalmente: preço (R\$ 50,00 a R\$ 350,00), notas olfativas (pirâmide 5-7 componentes), e marca.

\subsection{Matriz de Utilidade}

A matriz de utilidade foi construída sinteticamente com características realistas:
\begin{itemize}
    \item \textbf{Dimensões:} 500 usuários × 50 produtos = 25.000 células potenciais
    \item \textbf{Avaliações:} 5.905 registradas (23,6\% de densidade)
    \item \textbf{Escala:} 1-5 (inteiros), distribuição realista por persona de usuário
    \item \textbf{Densidade:} Adequada para SVD (recomenda-se $>$ 10\%)
\end{itemize}

A geração de dados é reproduzível e determinística (seed=42): qualquer execução posterior regenera a mesma matriz de utilidade, garantindo validação consistente e comparabilidade dos resultados.

\section{Metodologia}

\subsection{Fase 1: Filtragem por Conteúdo (TF-IDF)}

O algoritmo TF-IDF gera candidatos iniciais baseado no perfil declarado do usuário.

\subsubsection{Fundamentos Matemáticos}

Dado um documento $d$ e termo $t$, o score TF-IDF é:

\begin{equation}
\text{TF-IDF}(t, d) = \text{TF}(t, d) \times \text{IDF}(t)
\end{equation}

onde:
\begin{align}
\text{TF}(t, d) &= \log(1 + f_{t,d}) \quad \text{(frequência do termo)} \\
\text{IDF}(t) &= \log\left(\frac{N}{n_t}\right) \quad \text{(peso inverso global)}
\end{align}

com $N$ = 50 (total de perfumes) e $n_t$ = número de perfumes contendo termo $t$.

\textbf{Intuição:} Termos raros aumentam sua relevância discriminativa. Por exemplo, ``amadeirado noturno'' é mais específico que apenas ``amadeirado'', oferecendo maior poder de discriminação entre produtos similares.

\subsubsection{Fluxo de Vetorização}

\begin{enumerate}
    \item Concatenar características: ``floral'' + ``noturno casual'' + ``feminino''
    \item Tokenizar em scikit-learn (TfidfVectorizer)
    \item Gerar matriz TF-IDF 50×D (D = vocab size, pré-calculado)
    \item Calcular similaridade por cosseno: $\cos(\vec{u}, \vec{v}) = \frac{\vec{u} \cdot \vec{v}}{|\vec{u}||\vec{v}|}$
\end{enumerate}

\textbf{Resultado:} Score TF-IDF normalizado em [0, 1] para cada perfume. Aplicam-se filtros (gênero, preço), selecionam-se top-20.

\subsection{Fase 2: Filtragem Colaborativa (SVD)}

O SVD reordena os top-20 candidatos de TF-IDF usando padrões de usuários similares.

\subsubsection{Matrix Factorization}

A matriz de utilidade $R \in \mathbb{R}^{500 \times 50}$ é decomposta:

\begin{equation}
R = U \Sigma V^T \approx \hat{U} \hat{\Sigma} \hat{V}^T
\end{equation}

mantendo $k=50$ fatores singulares. As matrizes resultantes são:

\begin{align}
U \in \mathbb{R}^{500 \times 50} &\quad \text{(preferências latentes por usuário)} \\
V^T \in \mathbb{R}^{50 \times 50} &\quad \text{(características latentes por perfume)}
\end{align}

Multiplicação $U \Sigma V^T$ prediz todas as 25.000 células, inclusive as vazias.

\subsubsection{Reordenação Colaborativa}

\begin{enumerate}
    \item Identificar 50 usuários mais similares ao novo usuário (via TF-IDF no perfil)
    \item Extrair predições SVD para os top-20 perfumes: $\hat{R}[50 \text{ usuários}, \text{top-20}]$
    \item Calcular score médio: $s_j = \frac{1}{50} \sum_{i=1}^{50} \hat{R}_{ij}$
    \item Normalizar em [0, 1]: $s\_norm_j = \frac{s_j - 1}{4}$ (escala de notas 1-5)
\end{enumerate}

\subsection{Hibridização}

O score final combina ambas as abordagens:

\begin{equation}
\text{score}_{\text{final}}(j) = 0.7 \times \text{tfidf}(j) + 0.3 \times \text{svd\_norm}(j)
\end{equation}

\textbf{Justificativa dos pesos:}

\begin{itemize}
    \item \textbf{70\% TF-IDF:} Com apenas 50 produtos, conteúdo é discriminativo. Oferece explicabilidade alta (fácil justificar ao usuário).
    \item \textbf{30\% SVD:} Com 23,6\% densidade, captura padrões reais sem overfitting. Permite serendipidade (descoberta).
    \item \textbf{Balanço:} 70/30 oferece cobertura robusta sem sacrificar relevância.
\end{itemize}

\textbf{Threshold de qualidade:} Score $\geq 0.50$ retorna recomendação. Abaixo disso, ativa fallback.

\section{Cascata de Fallback}

Quando o modelo não encontra match satisfatório, implementa-se 4 níveis progressivos:

\begin{description}
    \item[Nível 0 (Match Perfeito)] Todas as restrições (família + ocasião + gênero + preço), com validação que família aparece no produto. Mensagem: ``✅ Recomendação personalizada''.
    
    \item[Nível 1 (Ocasião Restritiva)] Remove ocasião (algumas são raras). Score $\geq 0.50$. Mensagem: ``⏰ Poucas opções para essa ocasião''.
    
    \item[Nível 2 (Gênero Cruzado)] Remove gênero (famílias desbalanceadas por gênero). Score $\geq 0.50$. Mensagem: ``✨ Seu perfil é raro. Mostrando similares''.
    
    \item[Nível 3 (Top-5 Populares)] Sem filtros, retorna mais vendidos por gênero. Mensagem: ``Sem match ideal. Mostrando os perfumes mais populares.''
\end{description}

\section{Resultados e Validação}

O sistema foi testado com 100 cenários aleatórios:
\begin{itemize}
    \item Nível 0: 55\% (match perfeito muito comum)
    \item Nível 1: 25\% (ocasião é realmente restritiva)
    \item Nível 2: 15\% (gênero desbalanceado para algumas famílias)
    \item Nível 3: 5\% (combinações impossíveis raras)
\end{itemize}

\textbf{Cobertura:} 55\% das buscas encontram match satisfatório em Nível 0 (score $\geq 0.50$). Adicionando Níveis 1-2, \textbf{95\% alcançam recomendações de qualidade}. Apenas 5\% requerem fallback para Nível 3 (top-5 
populares).

\section{Conclusão}

O projeto implementa com sucesso um sistema híbrido que combina TF-IDF (explicabilidade + cold-start) e SVD (colaboração + serendipidade) na proporção 70/30. A cascata de fallback com 4 níveis garante cobertura (sempre retorna algo) sem sacrificar relevância (threshold 0.50).

\textbf{Contribuições principais:}
\begin{enumerate}
    \item Hibridização efetiva que resolve cold-start mantendo interpretabilidade
    \item Validação explícita de restrições (família, gênero, preço) reduzindo falsos positivos
    \item Cascata robusta de fallback com relaxamento inteligente de restrições
\end{enumerate}

A abordagem escala bem para catálogos pequenos-médios (50 produtos, 500 usuários) e pode ser estendida com retrainamento periódico de SVD conforme novos ratings acumulem.

\section*{Referências}

\begin{itemize}
    \item Ricci, F., Rokach, L., \& Shapira, B. (Eds.). (2015). \textit{Recommender Systems Handbook}. Springer.
    
    \item Koren, Y., Bell, R., \& Volinsky, C. (2009). Matrix factorization techniques for recommender systems. \textit{IEEE Computer}, 42(8), 30--37.
    
    \item Lops, P., De Gemmis, M., \& Semeraro, G. (2011). Content-based recommender systems: State of the art and trends. In \textit{Recommender systems handbook} (pp. 73--105). Springer.
\end{itemize}

\end{document}
